\section{Discussion: using SMCD in practice}
\label{sec:discussion}

\paragraph{A workflow.} For a new problem the method reduces to five steps, all but the
last done before any training.
\begin{enumerate}
  \item Write the operator's Fourier symbol and the forcing and boundary data
    (Proposition~2); for a nonlinear operator, linearise about a base state.
  \item Threshold the resulting spectrum to get the target frequencies $\hat S$, with
    band $K$ and spacing $\Delta$.
  \item Read off the depth $L=\lceil\log_3(2K/\Delta+1)\rceil$, the ternary scalings, the
    qubit count and the entangler (Algorithm~1), and check the design card: coverage
    should be 1 and the predicted benefit non-trivial.
  \item If the design card predicts little benefit, as it does for a smoothing operator,
    stop: a classical PINN is the better tool, and the answer cost no training.
  \item Otherwise build $K$ copies of the designed circuit behind a linear head and train
    it like any PINN; eight copies were enough for both problems here.
\end{enumerate}
The worked designs of Section~\ref{sec:worked} show the arithmetic is short enough to do
by hand, and the implementation does it in milliseconds.

\paragraph{Which problems suit the method.} SMCD pays off when the frequencies a
solution needs are known from the problem's structure before it is solved: evenly spaced
sources such as the canals here, wells or drains on a grid, periodic geometry, or forcing
at known harmonics. In those problems the classical difficulty is precisely spectral bias,
fine structure that a standard network learns last, and the design rule places the
circuit's frequencies exactly on that structure. The groundwater application shows the
practical consequence: the fine structure decides the planning answer (where the water
table crosses the root zone), and the circuit gets it right to 1\,m.

\paragraph{Explainability as a design check.} The instruments did more than report
accuracy: the per-frequency error showed that improvements land on the designed
frequencies, and the same instrument exposed why a single circuit trades a slow wave for a
fast one, which led directly to circuit copies. Running the same instruments at the same
checkpoints on every model turned them from after-the-fact illustrations into tests,
and each expected signature was stated before the runs.

\paragraph{Towards hardware.} Everything a PINN needs from the circuit is available on
quantum hardware without approximation. The circuit's output is an expectation value;
its derivatives with respect to the angles and to the input $x$ follow from the
parameter-shift rule (Section~\ref{sec:background}), so the residual $-u''-f$ and its
gradient can be estimated from repeated circuit evaluations alone. The circuits are as
short as a re-uploading model can be: one qubit, four encoding layers and five trainable
rotations per copy, with no entangling gates, which keeps them within the coherence times
of current devices. Eight copies can run as eight independent single-qubit circuits, on
separate qubits in parallel or one after another. The main cost on hardware is shot
noise in the estimated expectation values and their derivatives, which the noise
surrogates in the implementation are built to study.
